\capitulo{2}{Objetivos del proyecto}

En este capítulo se recogen los objetivos que orientan el desarrollo del proyecto. Se han organizado en cuatro bloques: el objetivo general, los objetivos funcionales (qué hace la aplicación), los objetivos no funcionales (cómo lo hace) y los objetivos personales que me planteo como autor.

\section{Objetivo general}

El objetivo general es desarrollar un prototipo \gls{standalone} en Python que reciba síntomas escritos en lenguaje natural y devuelva la especialidad médica más adecuada para gestionar la cita del paciente. Todo el procesamiento se ejecuta en local, sin enviar datos a ningún servicio externo, y la aplicación incluye un mecanismo de confianza con un umbral configurable: cuando el clasificador no está seguro de su predicción, el caso se marca para revisión en lugar de devolver un resultado dudoso.

Sobre esa base, el trabajo persigue cuatro metas más concretas:

\begin{itemize}
	\item Implementar un flujo de \gls{pln} completo, desde el preprocesado hasta la predicción, dentro de un único ejecutable.
	\item Comparar cuatro enfoques de clasificación sobre el mismo conjunto de datos y bajo el mismo protocolo: \gls{tfidf} como línea base, \gls{tfidf} con reducción \gls{svd}, un clasificador \gls{svm} y una \gls{rna} de tipo Perceptrón Multicapa (\gls{mlp}).
	\item Medir no solo el acierto de cada modelo, sino también su coste en tiempo y memoria, para identificar cuál ofrece el mejor equilibrio en un escenario con recursos limitados.
	\item Demostrar que un sistema autónomo y ligero puede servir como apoyo al triaje en entornos donde la conectividad no está garantizada o donde los datos del paciente no deben salir del equipo.
\end{itemize}

\section{Objetivos funcionales}

Los objetivos funcionales describen qué debe hacer la aplicación desde el punto de vista de quien la utiliza. El sistema atiende a dos perfiles diferentes (paciente y administrador).

\subsection{Preparación de los datos}

La aplicación parte de un CSV con tres columnas: descripción de síntomas, enfermedad asociada y especialidad médica que correspondería asignar. Antes de entrenar nada, los datos pasan por una fase de validación que detecta nulos, duplicados, etiquetas que no figuran en el catálogo y desequilibrios fuertes entre clases.

Una vez saneados los datos, se aplica un preprocesado que limpia el texto, lo divide en palabras y genera combinaciones de dos o tres palabras consecutivas para captar expresiones como ``dolor torácico'' o ``falta de aire''. Este mismo preprocesado se aplica después a cualquier síntoma que introduzca el paciente, de manera que el texto llegue al modelo en el mismo formato con el que fue entrenado.


\subsection{Clasificación con cuatro modelos}

El núcleo del sistema son cuatro clasificadores entrenados de forma paralela: \gls{tfidf} solo, \gls{tfidf} combinado con \gls{svd} mediante \gls{truncatedsvd}, un \gls{svm} y un \gls{mlp}. La idea no es elegir uno a priori sino dejar que sea la evaluación posterior la que diga cuál funciona mejor en este dominio. El \gls{mlp} se entrena con \textit{early stopping} y regularización para evitar sobreajuste \cite{ibm_redes_neuronales, aws_sobreajuste}.

\subsection{Predicción a partir de texto libre}

El paciente introduce sus síntomas (por escrito o por voz, gracias a un módulo de transcripción) y la aplicación devuelve la especialidad predicha junto con el porcentaje de confianza. Si esa confianza queda por debajo del umbral configurado, el sistema no se rinde: pide al paciente que amplíe la descripción o que escoja entre las dos o tres alternativas más probables. Si tras esa segunda vuelta la predicción sigue siendo dudosa, la consulta se deriva a atención general y queda registrada en un log para revisarla más adelante.

\subsection{Panel de administración}

Un segundo perfil, el administrador, tiene acceso a una pestaña separada con tres herramientas. La primera permite añadir nuevos pares síntoma-especialidad, que se guardan en un CSV ampliado distinto del original para no contaminar los datos de partida. La segunda muestra el log de consultas no clasificadas con confianza suficiente, donde el administrador puede asignar manualmente la especialidad correcta y volcar el caso al mismo CSV ampliado. La tercera dispara el reentrenamiento de los cuatro modelos con los datos actualizados, sin necesidad de tocar el código.

\subsection{Evaluación y comparación}

Los cuatro modelos se evalúan con precisión, \gls{recall} y F1 (macro y micro), y se acompañan de matrices de confusión que ayudan a identificar entre qué especialidades se concentran los errores. Como los textos clínicos reales pueden ser muy variables, también se comprueba qué ocurre cuando la entrada es corta, contiene faltas de ortografía o usa vocabulario coloquial en lugar del término médico.

\section{Objetivos no funcionales}

\subsection{Ejecución local}

La aplicación tiene que funcionar sin conexión a internet. Esta decisión condiciona muchas otras: obliga a elegir modelos ligeros que quepan en memoria y se ejecuten en CPU, a serializar todo en disco para no depender de servicios remotos y a renunciar a librerías o servicios que requieran una clave o un endpoint externo. Como efecto colateral, los datos del paciente nunca salen del equipo.

\subsection{Sin dependencia de inteligencia artificial externa}

El proyecto descarta de forma deliberada el uso de modelos de inteligencia artificial entrenados por terceros (LLMs comerciales, servicios en la nube de clasificación clínica, etc.) y tampoco contempla entrenar una IA propia de gran tamaño. Hay dos razones para esta decisión. La primera es económica: los LLMs comerciales tienen un coste por consulta que haría inviable su uso en un sistema sanitario público, y entrenar una red neuronal grande desde cero requiere una infraestructura que está fuera del alcance de un TFG. La segunda es operativa: cualquier servicio externo obliga a depender de internet y a enviar datos del paciente a servidores ajenos, lo que choca con el planteamiento de privacidad y autonomía del prototipo. Por eso la elección recae en técnicas estadísticas y de aprendizaje supervisado clásico, que pueden entrenar en minutos sobre un equipo modesto y se ejecutan sin depender de nadie.

\subsection{Modularidad}

El código se reparte entre módulos con responsabilidades bien delimitadas: carga de datos, preprocesado, vectorización, modelos, evaluación e interfaz. Esta separación no es solo cuestión de orden; permite, por ejemplo, sustituir la entrada de texto por una transcripción de llamadas telefónicas sin tener que modificar los clasificadores.

\subsection{Rendimiento razonable}

El proyecto no aspira a competir con el estado del arte en precisión. Lo que busca es encontrar el punto en el que un modelo lo bastante bueno se puede entrenar y ejecutar en un equipo modesto. Por eso se miden los tiempos de entrenamiento y de predicción y el consumo de memoria de cada uno de los cuatro modelos, y la decisión final se toma combinando esos datos con las métricas de acierto.

\section{Objetivos personales}

Más allá del producto final, hay varias cosas que quiero llevarme de este trabajo:

\begin{itemize}
	\item Llevar a la práctica lo aprendido sobre clasificación supervisada en un problema real. El CSV inicial no está balanceado (atención general y medicina interna concentran muchos más casos que especialidades como dermatología o neurología), pero el propio sistema está pensado para corregir ese desequilibrio: cada vez que el administrador revisa una consulta que no se clasificó con confianza y la añade al dataset ampliado, las especialidades menos representadas van ganando ejemplos y los modelos mejoran al reentrenarse.
	\item Entender, comparando modelos en igualdad de condiciones, cuándo compensa el tiempo de entrenamiento y la dificultad de ajustar una red neuronal y cuándo basta con un \gls{svm} sobre \gls{tfidf}, mucho más rápido y sencillo de afinar. Es una intuición que solo se adquiere implementando ambos enfoques sobre los mismos datos.
	\item Cerrar TFG con un producto funcional demostrable y una memoria técnica sólida, cumpliendo los requisitos académicos necesarios para finalizar el Grado en Ingeniería Informática.
\end{itemize}